
\section{The arity-one equation: solvability and free directions}
\label{sec:at-eq1}

\begin{theorem}[The arity-one equation]
\label{thm:at-eq1}
Let $F_1^{(0)}$ solve the arity-zero equation, $\Lambda^{(0)}F_1^{(0)}=\Ob^{(0)}$, and let $x$ be a generator.
\begin{itemize}
\item (i) $\mathcal R^{(1)}(x)$ is $\Lambda^{(0)}$-closed, of degree $0$ and Weyl weight $4-\tfrac12\lambda_x$, of position degree at most $p_x$, with the covariance of Theorem \ref{thm:at-bound}.
\item (ii) The equation $\Lambda^{(0)}F_1^{(1)}(x)=\mathcal R^{(1)}(x)$ has a local solution exactly when the class $[\mathcal R^{(1)}(x)]$ vanishes in the cohomology $H^0(\Lambda^{(0)})$ of the local functionals with these quantum numbers.
\item (iii) Two solutions differ by a $\Lambda^{(0)}$-closed local functional of degree $-1$ with the same quantum numbers; modulo $\Lambda^{(0)}$-exact ones, by a class of $H^{-1}(\Lambda^{(0)})$.
\item (iv) Replacing $F_1^{(0)}$ by $F_1^{(0)}+G^{(0)}$ with $\Lambda^{(0)}G^{(0)}=0$ replaces $\mathcal R^{(1)}(x)$ by $\mathcal R^{(1)}(x)-\{G^{(0)},S(x)\}$; for $G^{(0)}=\Lambda^{(0)}\phi^{(0)}$ the change is $+\Lambda^{(0)}\{\phi^{(0)},S(x)\}$, absorbed by $F_1^{(1)}(x)\mapsto F_1^{(1)}(x)+\{\phi^{(0)},S(x)\}$ (the field redefinition generated by $\phi^{(0)}$, Proposition c8-laws (a)).
\item (v) The map $x\mapsto[\mathcal R^{(1)}(x)]$ satisfies, in $H^0(\Lambda^{(0)})$,
\[
[\mathcal R^{(1)}([x,y])]=(-1)^{|x|}[\{\mathcal R^{(1)}(x),S(y)\}]-[\{S(x),\mathcal R^{(1)}(y)\}],
\]
the Chevalley--Eilenberg one-cocycle condition of $\mathfrak{psl}(4|4)$ with values in $H^0(\Lambda^{(0)})$ for the action through the brackets with $S(y)$; the change (iv), $\mathcal R^{(1)}(x)\mapsto\mathcal R^{(1)}(x)-\{G^{(0)},S(x)\}$, satisfies the same identity with the left side included (it is the coboundary of $[G^{(0)}]$). The obstruction to solving through arity one, for some choice of the arity-zero solution, is its class in $H^1\big(\mathfrak{psl}(4|4);H^0(\Lambda^{(0)})\big)$.
\end{itemize}
\end{theorem}

\begin{proof}
(i) Lemma \ref{lem:at-sformulas} (iii). (ii), (iii): definition of the cohomology; locality of the solution is part of the statement of $H^0$ on local functionals. (iv) $\{\Lambda^{(0)}\phi^{(0)},S(x)\}=-\Lambda^{(0)}\{\phi^{(0)},S(x)\}$ by the derivation rule of the proof of Lemma \ref{lem:at-sformulas} and $\Lambda^{(0)}S(x)=0$. (v) The arity-two identity of Lemma \ref{lem:at-sformulas} (ii) for $G=\Ob-sF_1$, whose arity-zero component vanishes, reads $\Lambda^{(0)}G^{(2)}(x,y)=G^{(1)}([x,y])-(-1)^{|x|}\{G^{(1)}(x),S(y)\}+\{S(x),G^{(1)}(y)\}$; in $H^0$ the left side vanishes, which is the stated identity with $G^{(1)}=\mathcal R^{(1)}$. For $\mathcal R^{(1)}(x)=-\{G^{(0)},S(x)\}$ the identity follows from $(-1)^{|x|}\{S(x),S(y)\}=-S([x,y])-\Lambda^{(0)}S^{(2)}(x,y)$ and the derivation rule $\{G^{(0)},\{a,b\}\}=-\{\{G^{(0)},a\},b\}-(-1)^{|a|}\{a,\{G^{(0)},b\}\}$, the term $\{G^{(0)},\Lambda^{(0)}S^{(2)}\}=-\Lambda^{(0)}\{G^{(0)},S^{(2)}\}$ being exact.
\end{proof}

\subsection{Two legs}

On two legs the bracket with $I$ raises the number of legs, so the equation reads $QF_1^{(1)}(x)|_2=\mathcal R^{(1)}(x)|_2$ and the relevant cohomology is that of $Q$ on quadratic functionals.

\begin{proposition}[Two legs along $\Sigma$]
\label{prop:at-xi2}
Let $\xi\in S_+\otimes W$ (for $S_-\otimes W^*$ exchange the chiralities). Write $\Gamma_V:=\Gamma(\xi,\lambda)|_V$, $\Gamma_\perp:=\Gamma(\xi,\lambda)|_{W^\perp_{\mathbb R}}$ and $(\xi,\lambda^*)$ for the $\mathfrak g$-valued pairings, and
\[
g_1=\int(\Gamma_V,d^*d\AV),\quad g_2=\int(\Gamma_V,dd^*\AV),\quad g_3=\int(\Gamma_\perp,d^*d\Ap),\quad g_4=\int\big(d^*dc,(\xi,\lambda^*)\big),\quad T_0=\int\big(d^*\AV,(\xi,\lambda^*)\big).
\]
\begin{itemize}
\item (i) The quadratic local functionals of degree $-1$, Weyl weight $\tfrac92$, linear in $\xi$ and covariant, are spanned by $T_0,T_1,T_2,T_3,T_4$; those of degree $0$ by $g_1,g_2,g_3,g_4$ (legs $(\lambda_-,\AV)$, $(\lambda_+,\Ap)$, $(c,\lambda^*_+)$).
\item (ii) $QT_1=-QT_3=q_1g_1$, $QT_2=-QT_4=q_3g_3$, $QT_0=a\,g_2+b\,g_4$ with $q_1,q_3,a,b$ nonzero; in the normalisation of Definition \ref{def:at-names} with the left-derivation rule for $Q=\{S,-\}$, both engine routes give $q_1=q_3=-\tfrac12$, $a=b=-1$. The kernel of $Q$ on the degree $-1$ space is spanned by $S^V_2(\xi)=T_1+T_3$ and $S^\perp_2(\xi)=T_2+T_4$; no nonzero combination is $Q$-exact.
\item (iii) The $Q$-closed functionals of degree $0$ are spanned by $g_1$, $g_3$ and $a\,g_2+b\,g_4$, which is the image of $Q$: the cohomology $H^0(Q)$ of these quadratic functionals is zero.
\item (iv) Consequently the two-leg part of the arity-one equation along $\xi$ is solvable for every $F_1^{(0)}$ solving arity zero: with $\mathcal R^{(1)}(\xi)|_2=r_1g_1+r_3g_3+r_0(ag_2+bg_4)$,
\[
F_1^{(1)}(\xi)|_2=\frac{r_1}{q_1}T_1+\frac{r_3}{q_3}T_2+r_0T_0+u\big(\theta^\xi_VS^V_2(\xi)+\theta^\xi_\perp S^\perp_2(\xi)\big),
\]
with free $\theta^\xi_V,\theta^\xi_\perp$. The coefficient $r_0$ of the gauge-fixing structure $T_0$ is fed by the finite counterterm $\tfrac{27}{16}u\int(d^*\AV,d^*\AV)$ through $\{F_1^{(0)},T_1\}\propto g_2$, and by the parts $g_2$, $g_4$ of $\Ob^{(1)}(\xi)|_2$.
\item (v) Wess--Zumino on two legs: $Q\,\Ob^{(1)}(\xi)|_2=-\{\Ob^{(0)}|_2,T_1\}$ with $\Ob^{(0)}|_2=-\tfrac{27}8u\int(d^*dc,d^*\AV)$ (Theorem c4-obstruction): the components of $\Ob^{(1)}(\xi)|_2$ along $g_2$, $g_4$ outside the closed combination are fixed by the arity-zero two-leg obstruction (gate G-WZ1).
\item (vi) The $\beta$-directions of $F_1^{(0)}$ (Theorem c9-F3) are $Q$-exact on two legs, $Q$ of the field-counting functionals; through $\{Q N_f,S^{(1)}_2(\xi)\}=-Q\{N_f,S^{(1)}_2(\xi)\}$ they shift $F_1^{(1)}(\xi)|_2$ by $+\{N_f,S^{(1)}_2(\xi)\}$ (Theorem \ref{thm:at-eq1} (iv)), which multiplies each arrow $f\to f'$ of the classical transformation by the difference of the counting weights (Proposition qs-laws (a)).
\end{itemize}
\end{proposition}

\begin{proof}
(i) $\xi\otimes\lambda_-\in(S_+\otimes S_-)\otimes(W\otimes W^*)$ contains the singlet of $SU(4)$ only in $V\otimes\mathbf 1$, so every bilinear in $(\xi,\lambda_-)$ is a component of $\Gamma_V$; with $\AV$ and two derivatives, after integration by parts onto $\AV$, the contractions are $\Gamma_V\cdot d^*d\AV$ and $\Gamma_V\cdot dd^*\AV$ (the $\epsilon$-contraction with two symmetric derivatives vanishes). $\xi\otimes\lambda_+\in(\mathbf 1\oplus\Lambda^2_+)\otimes(\mathbf 6\oplus\mathbf{10})$ pairs with $\Ap\in\mathbf 6$ through $\Gamma_\perp$, and the $\Lambda^2_+$ part needs an antisymmetric pair of derivatives; $(\xi\otimes\lambda_+,\AV)$ and $(\xi\otimes\lambda_-,\Ap)$ have no $SU(4)$-singlet. $\xi\otimes\lambda^*_+$ contains the scalar $(\xi,\lambda^*)$, paired with $c$ and two derivatives; $\xi\otimes\lambda^*_-$ has no singlet. Degree $-1$: $(\lambda,A^*)$ without derivatives gives $T_1$, $T_2$; $(A,\lambda^*)$ with one derivative gives $(\xi,\lambda^*_+)\,d^*\AV$, the self-dual part of $d\AV$ against $\Lambda^2_+$ ($T_3$), and $\rho(d\Ap)\xi$ against $\lambda^*_-$ ($T_4$); the weight count of Theorem \ref{thm:at-bound} excludes the rest. (ii) $Q$ substitutes $A^*\mapsto-\tfrac12d^*dA$, $\lambda^*\mapsto\dslash\lambda$, $\AV\mapsto-dc$ (Proposition c1-Q): $T_1\mapsto\Gamma_V\cdot d^*d\AV$, $T_2\mapsto\Gamma_\perp\cdot d^*d\Ap$; in $T_0$ the substitution in $\AV$ gives $g_4$ and that in $\lambda^*$ gives $(\xi,\dslash\lambda)\,d^*\AV$ with $(\xi,\dslash\lambda)=\sum_\mu(e_\mu,\Gamma(\xi,\partial_\mu\lambda))=-d^*\Gamma_V$, hence $g_2$; $S^V_2$, $S^\perp_2$ are closed as the two sector parts of the quadratic part of $\{S,S(\xi)\}=0$ (Proposition qs-freexi). The images $g_1$, $g_3$, $ag_2+bg_4$ are independent, so the kernel is two-dimensional. A $Q$-exact functional of degree $-1$ would be $Q$ of a degree $-2$ quadratic functional of weight $\tfrac92$, which does not exist (Proposition qs-freexi (i)). (iii) $Qg_1=-\int(\Gamma_V,d^*d\,dc)=0$, $Qg_3=0$ ($\Ap$ and $\lambda$ are inert), and $Qg_2=-\int(\Gamma_V,dd^*dc)$, $Qg_4=\int(d^*dc,(\xi,\dslash\lambda))=-\int(dd^*dc,\Gamma_V)$ are proportional and nonzero; the closed space is three-dimensional and contains the three-dimensional image. (iv) Lemma \ref{lem:at-sformulas} (iii) makes $\mathcal R^{(1)}(\xi)|_2$ closed; (iii) and (ii). The bracket $\{\int(d^*\AV,d^*\AV),T_1\}$ contracts $\AV$ with $A^*|_V$ and equals $\pm2\int(d^*\Gamma_V,d^*\AV)=\pm2g_2$. (v) Lemma \ref{lem:at-sformulas} (ii) on two legs. (vi) The quadratic part of $s$ of $N_f$ is $QN_f$, and $\{QN_f,S_2\}=-Q\{N_f,S_2\}$ by the derivation rule and $QS_2=0$.
\end{proof}

\begin{proposition}[Two legs along the other generators]
\label{prop:at-two}
\leavevmode
\begin{itemize}
\item (i) $\Delta$. The $Q$-closed quadratic functionals of degree $-1$, weight $4$, position degree at most one and reflection-even are spanned modulo $Q$-exact ones by the sector parts $S(\Delta)_\perp,S(\Delta)_\lambda,S(\Delta)_V$, the exact ones by $Q\int(\iota_x\AV,c^*)$ and $Q\int(\rho(x)\lambda^*,\lambda^*)$ (Proposition c11-tc0). The two-leg equation is solved by the family of Definition c11-free once $\Ob^{(1)}(\Delta)|_2$ is known; the strict route computes $\Ob^{(1)}(\Delta)|_2$ by \eqref{eq:at-localformula}, regression target
\begin{equation}
\hbar\Ob^{(1)}(\Delta)\big|_2=u\Big[-\tfrac{11}3\int(dA,dA)+\tfrac{47}3\int(\lambda,\dslash\lambda)-\tfrac{14}3\int(dc,A^*)\Big]
\label{eq:at-reg2}
\end{equation}
(Theorem c9-quadratic (ii)), with coefficient $0$ for $\int(d^*\AV,d^*\AV)$ and for every functional with an explicit factor $x$; by Proposition \ref{prop:at-pos} the last statement is $\Ob^{(1)}(P_b)|_2=0$.
\item (ii) $P_b$. The quadratic $Q$-closed functionals of degree $-1$, weight $5$, linear in $b$ and translation invariant are spanned modulo exact ones by $S(P_b)_\perp,S(P_b)_\lambda,S(P_b)_V$ (the position parts of those of (i), Proposition \ref{prop:at-pos}); the exact ones by $Q\int(\iota_b\AV,c^*)$, $Q\int(\rho(b)\lambda^*,\lambda^*)$. Completeness at this weight is task S1-P.
\item (iii) $\mathfrak{so}(V)$, $\mathfrak{so}(W^\perp_{\mathbb R})$. The sector parts of $S(M_\omega)$, of $S(R_r)$ ($\perp$ and $\lambda$ only) are closed; along $R_r$, $F_1^{(1)}(R_r)=0$ solves the equation (Theorem \ref{thm:at-second} (iii)). Along $M_\omega$, in the leg sets $(\AV,\AV)_2$, $(c,A^*|_V)_1$ the closed functionals of degree $0$ are spanned by $Q\int\omega_{\mu\nu}(A_\mu,A^*_\nu)$, so $H^0(Q)=0$ there, and the spin term $\int\omega_{\mu\nu}(A_\mu,A^*_\nu)$ of $F_1^{(1)}(M_\omega)$ is fixed by $\Ob^{(1)}(M_\omega)|_2$. In the leg set $(\lambda_+,\lambda_-)_1$ the $\omega$-linear functionals of degree $0$ are spanned by $Q\int(\tfrac14\omega_{\mu\nu}\rho_\mu\rho_\nu\lambda,\lambda^*)$, the structure $\int(\lambda,\omega_{\mu\nu}\rho_{[\mu}\rho_\nu\rho_{\rho]}\partial_\rho\lambda)$ being a total derivative ($\Gamma$ symmetric), so $H^0(Q)=0$ there as well; with $S(M_\omega)_{\rm spin}:=\int\omega_{\mu\nu}(A_\mu,A^*_\nu)+\int(\tfrac14\omega_{\mu\nu}\rho_\mu\rho_\nu\lambda,\lambda^*)$ both engine routes give $\hbar\Ob^{(1)}(M_\omega)|_2=\tfrac13u\,Q\,S(M_\omega)_{\rm spin}$, uniformly in the two sectors, and $0$ on $(\Ap,\Ap)$ (Section \ref{sec:at-computed}).
\item (iv) $\eta\in S_+\otimes W^*$. The translation-invariant quadratic functionals of degree $-1$ and weight $\tfrac72$ are spanned by $\int(\rho(\AV)\eta,\lambda^*)$ and $\int s_0$; those of degree $0$ by $\int(\eta,\lambda_+)\,d^*\AV$, $\int(\eta,\rho(d\AV)\lambda_+)$, the structure $(\lambda_-,\Ap)_1$ and the structure $(c,\lambda^*_-)_1$. $Q$ is block triangular for the split into position parts (Proposition \ref{prop:at-pos}) and intrinsic parts; its kernel on the full degree $-1$ space is spanned by $S^V(\eta)$ and $S^\perp(\eta)$ (route B of Part Q, task TQ0 along $\Sigma^*$). The cohomology $H^0(Q)$ on the full degree-$0$ two-leg space at weight $\tfrac72$ with position degree at most one (task S1-$\eta$) is one-dimensional, a class in the $(\Ap,\lambda)$ sector (both routes); the computed residual $\mathcal R^{(1)}(\eta)|_2$ has no component along it, so the two-leg equation along $\eta$ is solvable, with the two free directions $S^V(\eta)$, $S^\perp(\eta)$ (Section \ref{sec:at-computed}).
\item (v) $K_b$. The sector parts $S(K_b)_d$ are closed; the intrinsic two-leg target has the leg sets $(c,A^*|_V)_0$, $(\AV,\AV)_1$, $(\Ap,\Ap)_1$, $(\lambda_+,\lambda_-)_0$ (Table \ref{tab:at-census1}); $H^{-1}(Q)$ and $H^0(Q)$ with position degree at most two (task S1-K) have dimensions $7$ and $4$ (both routes): $H^{-1}$ is spanned by the three $[S(K_b)_d]$ and four second-order free-field symmetries; the residual has no component on $H^0$ (Section \ref{sec:at-computed}). $\{\Ob^{(0)}|_2,S(K_b)\}\neq0$ because $\int(d^*dc,d^*\AV)$ is not invariant under the special conformal transformations, and $\{F_1^{(0)},S(K_b)\}\neq0$ for the same reason.
\end{itemize}
\end{proposition}

\begin{proof}
(i) Proposition c11-tc0 and Theorem c9-quadratic (iii); the regression values are stated in Theorem c9-quadratic (ii). (ii) The position part along $b$ commutes with $Q$ (Lemma c11-canonical) and maps the closed and exact directions of (i) onto those listed. (iii) $v.$ commutes with $Q$; the leg-set computation in the vector sector is: $Q\int\omega_{\mu\nu}(\partial_\mu A_\nu,d^*\AV)$ and $Q\int\omega_{\mu\nu}(\partial_\mu c,A^*_\nu)$ are proportional and nonzero, so the closed subspace of these two is one-dimensional and contains $Q\int\omega_{\mu\nu}(A_\mu,A^*_\nu)\ne0$. (iv) The census of Table \ref{tab:at-census1}: $\eta\otimes\lambda^*_-\in V\otimes(\mathbf1\oplus\mathbf{15})$ pairs with $\AV$, $\eta\otimes\lambda^*_+\in(\mathbf1\oplus\Lambda^2_+)\otimes(\mathbf6\oplus\overline{\mathbf{10}})$ with $\Ap$; the degree-$0$ structures likewise. $Q$ commutes with $-\{S(P_b),-\}$, which gives the triangular form. (v) Census; $\{\int(d^*dc,d^*\AV),S(K_b)\}$ is the variation of a functional that is not conformally invariant.
\end{proof}

\subsection{Three and four legs: the sector coefficients and the arity-zero couplings}

\begin{theorem}[The cubic part fixes the sector differences]
\label{thm:at-cubic}
Let $x$ be a generator and write the quadratic part of a solution as $F_1^{(1)}(x)|_2=F^{\rm part}_2(x)+u\sum_d\theta^x_dS(x)_d+(\text{exact})$, with the sectors $d$ of Definition \ref{def:at-names} (ii). The part with three legs of the arity-one equation is
\[
Q\,F_1^{(1)}(x)|_3+\{I_3,F_1^{(1)}(x)|_2\}=\mathcal R^{(1)}(x)|_3,
\]
$I_3$ the cubic part of $I$, and $\mathcal R^{(1)}(x)|_3=\Ob^{(1)}(x)|_3-\{F_1^{(0)}|_2,S_3(x)\}-\{F_1^{(0)}|_3,S_2(x)\}$.
\begin{itemize}
\item (i) Each $\{I_3,S(x)_d\}$ is $Q$-closed, and $\sum_d\{I_3,S(x)_d\}=\{I_3,S_2(x)\}=-QS_3(x)$; the equation depends on the $\theta^x_d$ through $\sum_{d\ne d_0}(\theta^x_d-\theta^x_{d_0})[\{I_3,S(x)_d\}]$ in the cohomology $H^0(Q)_3$ of the cubic functionals of degree $0$, weight $4-\tfrac12\lambda_x$ and position degree at most $p_x$.
\item (ii) The cubic equation is solvable exactly when $[\mathcal R^{(1)}(x)|_3-\{I_3,F^{\rm part}_2(x)\}]=u\sum_{d\ne d_0}(\theta^x_d-\theta^x_{d_0})[\{I_3,S(x)_d\}]$ in $H^0(Q)_3$.
\item (iii) If the classes $[\{I_3,S(x)_d\}]$, $d\ne d_0$, are linearly independent, the cubic equation fixes the differences $\theta^x_d-\theta^x_{d_0}$. They are independent exactly when no combination $\sum_dc_dS(x)_d$ other than the multiples of $S_2(x)$ extends to a $\Lambda^{(0)}$-closed functional through three legs, that is, when the sector parts are not symmetries of the cubic interacting theory.
\item (iv) Along $\Delta$ this is the test GC13 of the composite workflow: the rank of $[\{I_3,S(\Delta)_\perp\}]$, $[\{I_3,S(\Delta)_\lambda\}]$ modulo $Q$ of the cubic functionals of degree $-1$, weight $4$ and position degree at most one; its position part along $b$ is the same test along $P_b$ with $S(P_b)_d$, and $\{F_1^{(0)},S(\Delta)\}=0$ makes $\mathcal R^{(1)}(\Delta)=\Ob^{(1)}(\Delta)$. Along $\xi$ the single class $[\{I_3,S^V_2(\xi)\}]$ decides $\theta^\xi_V-\theta^\xi_\perp$; along $\eta$, $[\{I_3,S^V(\eta)\}]$; along $K_b$, $\mathfrak{so}(V)$, $\mathfrak{so}(W^\perp_{\mathbb R})$, the classes of the brackets with the sector parts.
\end{itemize}
\end{theorem}

\begin{proof}
On three legs, $\Lambda^{(0)}$ applied to a functional with two and three legs gives $Q$ of the cubic part plus the bracket of $I_3$ with the quadratic part; $\{F_1^{(0)},S(x)\}$ has on three legs the two stated terms. (i) $QI_3=0$ (cubic part of the classical master equation, Proposition c9-maps (i)) and $QS(x)_d=0$ give $Q\{I_3,S(x)_d\}=0$ by the derivation rule; the cubic part of $\{S,S(x)\}=0$ is $QS_3(x)+\{I_3,S_2(x)\}=0$. (ii) Definition of $H^0(Q)_3$; the exact quadratic directions contribute $\{I_3,Q\psi\}=\pm Q\{I_3,\psi\}$, exact. (iii) A relation $\sum_dc_d[\{I_3,S(x)_d\}]=0$ means $\{I_3,\sum_dc_dS(x)_d\}=-QY_3$ for a cubic $Y_3$, i.e. $\Lambda^{(0)}(\sum_dc_dS(x)_d+Y_3)=0$ through three legs; the combination $\sum_dS(x)_d$ always extends, by $Y_3=S_3(x)$. (iv) Lemma c9-SDelta (i): $\{F,S(\Delta)\}=-v.F=0$ for $F$ of weight $4$; Proposition \ref{prop:at-pos}.
\end{proof}

\begin{proposition}[The arity-one equations along $\Sigma$ constrain the arity-zero couplings]
\label{prop:at-couplings}
Let $\mathcal H_4:=H^0(\Lambda^{(0)})$ on colour-singlet local functionals of degree $0$ and weight $4$, invariant under $\mathrm{Spin}(4)\times\mathrm{Spin}(6)$ and the reflection: the classical couplings modulo field redefinitions, containing the class of $S$ (the coupling direction of Theorem c8-obstruction (iv)) and the classes $\kappa_1,\dots,\kappa_4$ of Theorem c9-F3 that extend to $\Lambda^{(0)}$-closed functionals. For $\xi\in\Sigma$ let $\delta_\xi:\mathcal H_4\to H^0(\Lambda^{(0)})_{9/2}$, $[G]\mapsto[\{G,S(\xi)\}]$. Then:
\begin{itemize}
\item (i) the arity-one equations along $\Sigma$ are solvable for some arity-zero solution exactly when $[\mathcal R^{(1)}(\xi)]$ lies in the image of $\delta_\xi$ for every $\xi$, the same $[G]$ serving all $\xi$;
\item (ii) the arity-zero freedom left after the arity-one equations along $\Sigma$ is $\bigcap_\xi\ker\delta_\xi$, the supersymmetric classes; it contains $[S]$, since $\{S,S(\xi)\}=0$;
\item (iii) along $\Delta$, $P_b$, $\mathfrak{so}(V)$, $\mathfrak{so}(W^\perp_{\mathbb R})$ and $K_b$ the map $[G]\mapsto[\{G,S(x)\}]$ vanishes on $\mathcal H_4$: the classes have weight $4$, are translation and rotation invariant, and the gauge-invariant functionals of weight $4$ of the fields $A$, $\lambda$ are conformally invariant; these equations impose no condition on the arity-zero couplings.
\end{itemize}
In particular the three- and four-leg parts of the arity-one equations along $\Sigma$, with the cubic finite counterterm of Theorem c9-F3, fix the coefficients $\kappa_k$ and the $\beta$-completions up to the coupling direction and field redefinitions.
\end{proposition}

\begin{proof}
Theorem \ref{thm:at-eq1} (iv), (v). (iii): $\{G,S(x)\}=\mp v.G$ for even $x$ (Lemma c9-SDelta (i), Definition c7-vdot), zero for $G$ of weight $4$ and invariant under translations and rotations; for $K_b$, the functionals $\int(F_A,F_A)$, $\int(d_A\iota_wA,d_A\iota_wA)$, $\int(\lambda,\dslash_A\lambda)$, the Yukawa and quartic terms are conformally invariant with the weights $1$, $\tfrac32$ (Proposition c2-coeff), and the antifield terms of weight $4$ are invariant because $S$ is.
\end{proof}

\begin{hypothesis}[Classical symmetries]
\label{hyp:at-H}
For every generator $x$ (resp. pair $(x,y)$), every $\Lambda^{(0)}$-closed colour-singlet local functional of degree $-1$, of Weyl weight $4-\tfrac12\lambda_x$ (resp. $4-\tfrac12(\lambda_x+\lambda_y)$), of position degree at most $p_x$ (resp. $p_x+p_y$), linear in the Killing data and equivariant as in Theorem \ref{thm:at-bound}, is a linear combination of the functionals $S(z)$, $z\in\mathfrak{psl}(4|4)$ with these quantum numbers, plus a $\Lambda^{(0)}$-exact functional.
\end{hypothesis}

\begin{remark}[Status of the hypothesis]
\label{rem:at-H}
Hypothesis \ref{hyp:at-H} states that the global symmetries of the interacting classical theory at these weights are the superconformal ones. On two legs it fails: $H^{-1}(Q)$ contains the sector parts $S(z)_d$ and, at weight $5$ with position degree one (the weight of the pair $(\Delta,P_b)$), the second-order symmetries of the free fields built from the conformal Killing tensors $b^{(\mu}x^{\nu)}$. The hypothesis is the statement that none of these extends through three legs except the full $S(z)$; Theorem \ref{thm:at-cubic} (iii) is its test at the weights of the arity-one components, and task S1-cubic of the engines computes it, at the weights of the arity-two pairs as well.
\end{remark}

\begin{theorem}[Solution space at arity one]
\label{thm:at-H}
Assume Hypothesis \ref{hyp:at-H} and that the arity-one equations are solvable. Then for every generator $x$ the solutions are $F_1^{(1)}(x)=F^{\rm part}(x)+t_x\,u\,S(x)+\Lambda^{(0)}(\text{local})$ with one number $t_x$, the uniform coefficient: all sector coefficients $\theta^x_d$ are equal to $t_x$ up to the fixed differences of Theorem \ref{thm:at-cubic} (iii) contained in $F^{\rm part}$. In a scheme equivariant for translations, rotations and $\mathfrak{so}(W^\perp_{\mathbb R})$, $t_x$ depends only on the type of $x$: $t_\Delta,t_P,t_M,t_R,t_{\xi_+},t_{\xi_-},t_{\eta_+},t_{\eta_-},t_K$.
\end{theorem}

\begin{proof}
Theorem \ref{thm:at-eq1} (iii) and Hypothesis \ref{hyp:at-H}; the only functional $S(z)$ with the weight, position degree and Killing data of $x$ that is equivariant is $S(x)$ itself. Equivariance under the compact groups makes $t_x$ constant on each irreducible block.
\end{proof}

\begin{remark}[The free directions and the equations that fix them]
\label{rem:at-freetable}
The differences of the sector coefficients are fixed by the cubic part of the arity-one equation (Theorem \ref{thm:at-cubic}) and, along $\Sigma$, $\Sigma^*$, $P_b$, $\Delta$, $\mathfrak{so}(V)$, $\mathfrak{so}(W^\perp_{\mathbb R})$, already by the two-leg parts of the arity-two equations (computed, Section \ref{sec:at-computed}); the uniform parts by the arity-two equations of Section \ref{sec:at-ar2}:
\[
\begin{array}{l|l|l}
x & \text{closed quadratic directions} & \text{uniform part fixed by}\\
\hline
\Delta & S(\Delta)_\perp,S(\Delta)_\lambda,S(\Delta)_V;\ Q\int(\iota_x\AV,c^*),\ Q\int(\rho(x)\lambda^*,\lambda^*) & (\Delta,P_b),\ (\Delta,\xi),\ (\xi,\eta)\\
P_b & S(P_b)_d;\ Q\int(\iota_b\AV,c^*),\ Q\int(\rho(b)\lambda^*,\lambda^*) & \text{up to the inner direction}\\
\mathfrak{so}(V) & S(M_\omega)_d & (K_b,P_{b'}),\ (\xi,\eta)\\
\mathfrak{so}(W^\perp_{\mathbb R}) & S(R_r)_\perp,\ S(R_r)_\lambda & (\xi,\eta)\\
\Sigma & S^V_2(\xi),\ S^\perp_2(\xi);\ \text{difference fixed by }(\xi,\xi'),\ (\xi_\pm,\xi'_\pm) & \text{up to inner and outer}\\
\Sigma^* & S^V(\eta),\ S^\perp(\eta);\ \text{difference fixed by }(\xi,\eta_\mp) & \text{up to inner and outer}\\
K_b & S(K_b)_d;\ \text{exact directions: task S1-K} & \text{up to the inner direction}
\end{array}
\]
\end{remark}
